% 来源：OEBPS/Text/01_54.xhtml（OPF spine 顺序；仅第一册）
\cleardoublepage
\BookEntry[安迪·格鲁夫《只有偏执狂才能生存》]{安迪·格鲁夫《只有偏执狂才能生存》}

\BookCover{images/书单37只有偏执狂才能生存.jpg}{[美] 安迪·格鲁夫著| 安然、张万伟译| 中信出版社2013年版}

\begin{BookSummaryBox}
\textbf{英特尔资深顾问安迪·格鲁夫的战略转折点法则}
\par
\end{BookSummaryBox}

\subsubsection*{这本书谈什么？}

本书是英特尔公司的前CEO安迪·格鲁夫，以在Intel公司多年的管理经验和在斯坦福大学商学院研究生院讲授企业战略管理课程为基础，撰写的英特尔公司发展壮大的回忆录。他在书中写道，你的企业迟早会走到一个战略转折点，企业的根基瞬间发生剧变，技术、规则、竞争环境、行业形态等等，一切的一切都瞬间改变了。如果放任自流，那么战略转折点的破坏力足以摧毁一家好的企业。但从另一方面说，战略转折点并不总是通向灾难，有敏锐眼光的管理者可以通过这个时机使企业跃升到新的高度。
\par

在担任英特尔首席执行官的11年中，安迪·格鲁夫多次被推到战略转折点的边缘。但他对转折点的思考，却帮助英特尔在残酷的竞争中生存并发展下去。在格鲁夫的领导下，英特尔成为世界上最大的电脑芯片公司，并在美国《财富》杂志500家最赚钱的公司中排名第7位。
\par

\subsubsection*{芒格为什么谈起这本书？}

在斯坦福大学法学院的演讲中，有人问芒格如何创造出一种轻松的氛围，让人们放弃追求一贯性的倾向，并承认他们犯下的错误。比如，英特尔的某个人谈到他们的奔腾芯片遇到了问题，要从头再来，尽管在这么复杂的企业中，这么做是很困难的。这个人想问问芒格会怎么做。芒格解释说：
\par

\begin{quotation}
英特尔及其同行创造了一种协调的企业文化，便于各个团队解决前沿的科学问题。但伯克希尔不同，因为它是一家控股公司，权力相当分散，只有最重要的资本配置才由公司的高层来拍板。所以，伯克希尔基本上会选择我们非常欣赏的人来管理我们的附属公司。
\par
\end{quotation}

但芒格也表示，他不清楚自己或沃伦是否擅长安迪·格鲁夫的老本行。安迪·格鲁夫加入刚创办的英特尔时，是该公司的第四名员工。1997年担任董事会主席和CEO，现在仍是英特尔的资深顾问。
\par

\subsubsection*{《只有偏执狂才能生存》作者安迪·格鲁夫如是说}

“敢于并有能力面对现实才是出路。”
\par

“应该坚持把自己认为正确的事情（当然不能是非法的和违背基本伦理道德的）做下去，除了对得起自己的良心，你也会收获意外的惊喜。”
\par

“以客户需求为导向，以客户的根本利益为出发点，短期内企业可能会承受巨大的经营压力，但一旦能熬过去，企业收获的将是无坚不摧的战斗力、坚实的口碑和滚滚财源。”
\par
